\documentclass[border=5mm]{standalone}
% ==============================================================================
% SETUP.TEX - CONFIGURACIÓN GLOBAL DEL PROYECTO
% ==============================================================================
% Este archivo centraliza todos los paquetes y estilos.
% Debe incluirse en el preámbulo del libro principal y de los archivos standalone.
% \PassOptionsToPackage{gray}{xcolor}
% --- 1. CODIFICACIÓN E IDIOMA ---
% fix-cm habilita las fuentes Computer Modern en tamaños arbitrarios (por debajo
% de 5 pt, entre otros). Debe cargarse antes que fontenc.
\usepackage{fix-cm}
\usepackage[utf8]{inputenc}
\usepackage[T1]{fontenc}
\usepackage[spanish, es-tabla]{babel}  
\usepackage{csquotes} % Recomendado para babel y citas

\usepackage{import}
% mode=image: Usa el PDF pre-compilado, nunca intenta recompilar.
% Debes compilar cada figura manualmente antes de compilar el libro.
\usepackage[mode=image, subpreambles=false]{standalone}

% --- 2. MATEMÁTICAS Y CIENCIAS ---
\usepackage{amsmath, amssymb, amsfonts, mathtools}
\usepackage{enumitem}

% LaTeX trae \DeclareMathSizes{5}{5}{5}{5}, así que a 5 pt (\tiny) los subíndices
% salen del mismo tamaño que la base: en las etiquetas de figuras el M de
% $\omega_M$ queda desproporcionado. Se restablece la proporción habitual.
\DeclareMathSizes{5}{5}{3.7}{3}

% --- 3. DISEÑO Y GEOMETRÍA --- 
% Unificamos los márgenes para todo el documento
\usepackage[left=2.5cm,right=2.5cm,top=3cm,bottom=3cm]{geometry}
\makeatletter
\ifstandalone % Evita que geometry dañe el recorte de las figuras bajándolas 3cm
  \@ifclasswith{standalone}{crop=false}{
    % Es un capítulo/sección: Dejamos que geometry actúe.
  }{
    % Es una figura: Neutralizamos geometry para que no las recorte mal.
    \geometry{pass}
  }
\fi
\makeatother
\usepackage{parskip} % Espaciado entre párrafos sin sangría
\usepackage{float}   % Para usar [H] en figura

% Ajustes de flotantes para evitar espacios verticales exagerados
\renewcommand{\topfraction}{0.95}      % Permite que las figuras ocupen hasta el 95% superior
\renewcommand{\bottomfraction}{0.95}   % Permite que las figuras ocupen hasta el 95% inferior
\renewcommand{\textfraction}{0.05}     % Permite hojas con tan solo un 5% de texto
\renewcommand{\floatpagefraction}{0.8} % Requiere que una página de solo figuras esté 80% llena
\setcounter{topnumber}{4}
\setcounter{bottomnumber}{4}
\setcounter{totalnumber}{10}

% --- 4. GRÁFICOS Y TABLAS ---
\usepackage{graphicx}
\usepackage{booktabs} % Tablas profesionales
\usepackage{caption}  % Control de captions y \captionof
\usepackage{tikz}
\usetikzlibrary{arrows.meta, calc, angles, quotes, babel, backgrounds}
\usepackage{pgfplots}
\usepgfplotslibrary{groupplots}
\usepgfplotslibrary{external}
\usepgfplotslibrary{patchplots}
\pgfplotsset{compat=1.18}
\usepackage[american, straightvoltages]{circuitikz}

% --- 5. COLORES Y CAJAS ---

\usepackage[table]{xcolor}
\usepackage{tcolorbox}

% Definición de colores corporativos/estilo
\definecolor{utnblue}{RGB}{0, 51, 102}
\definecolor{codegreen}{rgb}{0,0.6,0}
\definecolor{codegray}{rgb}{0.5,0.5,0.5}
\definecolor{codepurple}{rgb}{0.58,0,0.82}
\definecolor{backcolour}{rgb}{0.96,0.96,0.96}

% --- 6. ENTORNOS PERSONALIZADOS (CAJAS) ---

% Caja "Resultado" (Azul Corporativo) — Definiciones, fórmulas clave, resultados importantes.
% Uso: \begin{resultado}[Fórmula de Euler] ... \end{resultado}
\newtcolorbox{resultado}[1][]{
    colback=utnblue!5!white,
    colframe=utnblue,
    title=\textbf{#1},
    fonttitle=\bfseries,
    sharp corners,
    boxrule=0.5mm
}

% Caja "Nota" (Gris) — Advertencias, aplicaciones, contexto secundario.
% Uso: \begin{nota}[Cuidado con la calculadora] ... \end{nota}
% Uso: \begin{nota} ... \end{nota}  → título por defecto "Nota"
\newtcolorbox{nota}[1][Nota]{
    colback=codegray!5!white,
    colframe=codegray,
    title=\textbf{#1},
    fonttitle=\bfseries,
    sharp corners,
    boxrule=0.5mm
}

% Caja inline para resaltar un valor y enlazarlo con su otra aparición en una tabla
% o en una cuenta: mismo color, mismo valor.
% Uso: \marcavalor{$4$}  o  \marcavalor[codegreen]{$4$}
\newtcbox{\marcavalor}[1][utnblue]{
    on line,
    boxsep=1pt, left=2pt, right=2pt, top=1pt, bottom=1pt,
    colback=#1!10!white,
    colframe=#1,
    boxrule=0.4pt,
    arc=2pt
}

% --- 7. CÓDIGO FUENTE (LISTINGS) ---
\usepackage{listings}

\lstdefinestyle{miestilo}{
    backgroundcolor=\color{backcolour},   
    commentstyle=\color{codegreen},
    keywordstyle=\color{magenta},
    numberstyle=\tiny\color{codegray},
    stringstyle=\color{codepurple},
    basicstyle=\ttfamily\footnotesize,
    breakatwhitespace=false,         
    breaklines=true,                 
    captionpos=b,                    
    keepspaces=true,                 
    numbers=left,                    
    numbersep=5pt,                  
    showspaces=false,                
    showstringspaces=false,
    showtabs=false,                  
    tabsize=2,
    frame=single,                    
    rulecolor=\color{codegray}
}

\lstset{style=miestilo}
% Soporte para tildes y ñ en bloques de código
\lstset{literate={ñ}{{\~n}}1 {á}{{\'a}}1 {é}{{\'e}}1 {í}{{\'i}}1 {ó}{{\'o}}1 {ú}{{\'u}}1}

% Operadores matemáticos personalizados

% Vector en sentido discreto (lista finita de coordenadas): negrita + flecha.
% Uso: \vect{v}, \vect{e}_k, \vect{0}.
\newcommand{\vect}[1]{\vec{\mathbf{#1}}}

% Fecha de versión y comando para capítulos standalone
\providecommand{\verdate}{\the\year.\ifnum\month<10 0\fi\the\month.\ifnum\day<10 0\fi\the\day}
\newcommand{\chapterversion}{%
    \vspace{-1.5cm}\begin{flushright}\small\textit{\color{codegray}Versión~\verdate}\end{flushright}\vspace{0.5cm}%
}

% --- 8. HIPERVÍNCULOS (DEBE SER EL ÚLTIMO PAQUETE) ---
\usepackage[colorlinks=true, linkcolor=utnblue, urlcolor=utnblue, citecolor=utnblue]{hyperref}

% --- 9. REFERENCIAS CRUZADAS ENTRE CAPÍTULOS STANDALONE ---
% Cuando se compila un capítulo standalone, xr-hyper lee los .aux de los
% otros capítulos (si existen) para resolver \ref{} a labels externos.
% En la compilación del libro completo este bloque no se activa.
\ifstandalone
  \usepackage{xr-hyper}
  \makeatletter
  \newcommand{\xrchap}[1]{%
    \edef\xrc@self{\detokenize{Capitulo#1}}%
    \edef\xrc@job{\jobname}%
    \ifx\xrc@self\xrc@job\else
      \IfFileExists{../#1capitulo/Capitulo#1.aux}%
        {\externaldocument{../#1capitulo/Capitulo#1}}{}%
    \fi
  }
  \makeatother
  \xrchap{01}\xrchap{02}\xrchap{03}\xrchap{04}
  \xrchap{05}\xrchap{06}\xrchap{07}\xrchap{08}
  \xrchap{09}\xrchap{10}\xrchap{11}\xrchap{12}
  \xrchap{13}
\fi

% Par de vectores giratorios conjugados c e^{j w_0 t} y conj(c) e^{-j w_0 t}
% cuya suma cae sobre el eje real. Imagen geometrica de la simetria
% conjugada (sec 5.3.6): la figura ilustra la cancelacion entre un
% vector rotando a +w_0 y otro a -w_0; el principio es identico para
% cualquier par (c_k, c_{-k}) de la serie con c_{-k} = conj(c_k).
%
% Ambos paneles dibujados con TikZ puro para que la escala vertical
% sea identica (1.3 cm/unidad). Panel (a) usa scale=1.3 uniforme;
% panel (b) usa x=5cm/unidad de tiempo, y=1.3cm/unidad de amplitud.

\begin{document}
\begin{tikzpicture}[>=Latex, font=\small]

% =====================================================
% Panel (a): plano complejo rotado 90 grados CCW
% Re vertical (hacia arriba), Im horizontal (positivo a la izquierda)
% =====================================================
\begin{scope}[scale=1.3]
    \draw[dashed, gray!50, thin] (0,0) circle (1);

    \draw[->, thick, black] (0, -2.2) -- (0, 2.2) node[right, font=\small] {Re};
    \draw[->, thick, black] (1.3, 0) -- (-1.3, 0) node[left, font=\small] {Im};

    \draw[thick, black] (-0.07, 2)  -- (0.07, 2);
    \draw[thick, black] (-0.07, -2) -- (0.07, -2);
    \node[anchor=east, font=\tiny] at (-0.10, 2)  {$2|c|$};
    \node[anchor=east, font=\tiny] at (-0.10, -2) {$-2|c|$};

    % Diagonales del paralelogramo (rombo) hacia la punta del sum
    \draw[dashed, gray!60, thin] ({-sin(45)}, {cos(45)}) -- (0, {2*cos(45)});
    \draw[dashed, gray!60, thin] ({ sin(45)}, {cos(45)}) -- (0, {2*cos(45)});

    % Vector suma sobre el eje Re
    \draw[->, thick, black] (0, 0) -- (0, {2*cos(45)});
    \fill[black] (0, {2*cos(45)}) circle (1pt);

    % Vector giratorio c e^{j w_0 t}: angulo 45 + 90 = 135 (rotado 90 CCW)
    \draw[->, thick, utnblue] (0, 0) -- ({-sin(45)}, {cos(45)});
    \fill[utnblue] ({-sin(45)}, {cos(45)}) circle (1pt);

    % Vector giratorio conj(c) e^{-j w_0 t}: angulo -45 + 90 = 45
    \draw[->, thick, red!80!black] (0, 0) -- ({sin(45)}, {cos(45)});
    \fill[red!80!black] ({sin(45)}, {cos(45)}) circle (1pt);

    % Flechas de sentido de rotacion
    \draw[->, thin, gray!85] ({1.13*cos(145)}, {1.13*sin(145)}) arc (145:168:1.13);
    \draw[->, thin, gray!85] ({1.13*cos(35)}, {1.13*sin(35)}) arc (35:12:1.13);

    \node[anchor=south east, font=\footnotesize, utnblue]
        at ({-sin(45)-0.04}, {cos(45)+0.08}) {$c\,e^{j\omega_0 t}$};
    \node[anchor=south west, font=\footnotesize, red!80!black]
        at ({sin(45)+0.04}, {cos(45)+0.08}) {$\bar{c}\,e^{-j\omega_0 t}$};
\end{scope}

% =====================================================
% Panel (b): coseno en el tiempo
% Mismo TikZ; x=5cm/unidad de t, y=1.3cm/unidad (igual que panel a)
% =====================================================
\begin{scope}[shift={(3.7cm, 0)}, x=5cm, y=1.3cm]
    \draw[->, thick, black] (0, -2.2) -- (0, 2.2);
    \draw[->, thick, black] (-0.04, 0) -- (1.15, 0) node[right, font=\small] {$t$};

    \draw[thick, black] ([xshift=-0.08cm]0, 2)  -- ([xshift=0.08cm]0, 2);
    \draw[thick, black] ([xshift=-0.08cm]0, -2) -- ([xshift=0.08cm]0, -2);
    \node[anchor=east, font=\tiny] at ([xshift=-0.115cm]0, 2)  {$2|c|$};
    \node[anchor=east, font=\tiny] at ([xshift=-0.115cm]0, -2) {$-2|c|$};

    \draw[thick, black] ([yshift=0.08cm]0.125, 0) -- ([yshift=-0.08cm]0.125, 0);
    \draw[thick, black] ([yshift=0.08cm]0.5, 0)   -- ([yshift=-0.08cm]0.5, 0);
    \draw[thick, black] ([yshift=0.08cm]1, 0)     -- ([yshift=-0.08cm]1, 0);
    \node[anchor=north, font=\tiny] at ([yshift=-0.10cm]0.125, 0) {$\tfrac{T_0}{8}$};
    \node[anchor=north, font=\tiny] at ([yshift=-0.10cm]0.5, 0)   {$\tfrac{T_0}{2}$};
    \node[anchor=north, font=\tiny] at ([yshift=-0.10cm]1, 0)     {$T_0$};

    % Coseno 2 cos(w_0 t); con t escalado a periodo 1 -> 2 cos(2 pi t)
    \draw[thick, black, smooth, samples=200, domain=0:1.05]
        plot (\x, {2*cos(deg(2*pi*\x))});

    % Vertical de la instantanea en t = T_0/8
    \draw[dashed, gray!70, thin] (0.125, 0) -- (0.125, 1.4142);
    \fill[black] (0.125, 1.4142) circle (1pt);
\end{scope}

% Linea punteada conectando la punta del vector suma (panel a) con
% el valor del coseno en t=T_0/8 (panel b). Ambos a TikZ y =
% 2cos(45) * 1.3 = 1.838 cm.
\draw[dotted, thick, gray!85]
    (0, {2*cos(45)*1.3}) -- ({3.7 + 0.125*5}, {2*cos(45)*1.3});

% Marcadores de panel (misma altura)
\node[anchor=south, font=\small] at (0, {2.2*1.3 + 0.3}) {(a)};
\node[anchor=south, font=\small] at ({3.7 + 0.555*5}, {2.2*1.3 + 0.3}) {(b)};

\end{tikzpicture}
\end{document}
